\begin{abstract}

\vspace{10pt}\par

Un encaje de $S^1$ en $S^2$ es una función $f:S^1\to S^2$ inyectiva y continua. Una función así es un homeomorfismo sobre su imagen, y por lo tanto $f(S^1)$ es una curva cerrada simple $\gamma$. ¿Qué sabemos sobre $\gamma$? Desde hace 150 años se sabe que $\gamma$ separa a $S^2$ en dos componentes conexas. Hace 100 años se probó que además existe un homeomorfismo $h:S^2\to S^2$ que manda $\gamma$ en $S^1$, lo que tiene como consecuencia que las componentes conexas de $\gamma^c$ tienen $\pi_1=0$. \par
\vspace{10pt}
\comm{qué sabemos sobre $ \gamma$ me parece muy general y no sé q cosas más "sabemos" sobre $\gamma$ q se me escapan}

¿Qué pasa en general con encajes de $f:S^{n-1}\to S^n$? Desde hace 100 años sabemos que $f(S^{n-1})^c$ tiene dos componentes conexas. Sin embargo, en estos casos no necesariamente existe un homeomorfismo $h:S^n\to S^n$ que mande $f(S^n)$ al encaje usual de $S^{n-1}\subset S^n$.\br

Estos dos útlimos resultados se probarán en esta monografía. Para el primero se introducirá al lector a la teoría de homología singular y se probarán algunos resultados sobre la misma. Para probar que el homeomorfismo $f$ no necesariamente se extiende a todo $S^n$, darán dos ejemplos de 2-esferas $S$ encajadas en $R^3$ con $\pi_1(S^c)\neq\varnothing$. Además veremos que los $\pi_1(S)$ serán ejemplos de grupos no triviales cuyo abelianizado es trivial.\br

Después veremos cositas de límites inversos y atractores de sistemas dinámicos.\comm{escribir xd}

\end{abstract}